% Transcription — fonds Alexandre Grothendieck, shelfmark 134-8, batch 4 (pages 61-80).
% Established from the facsimile https://grothendieck.umontpellier.fr/134-8.pdf, read on the
% private mirror archives/raw/134-8.pdf, per the skill transcribe-grothendieck.
% Pass: Opus 5.5 (claude-opus-5-5), 2026-10-03 — first pass, unchecked against the pages by a human.
%
% Limited pass for issue #46. Folder 134-8 belongs to Pursuing Stacks, edited by
% G. Maltsiniotis, so only what the edition does not carry is transcribed: the typed, signed
% instructions of p. 63 for copying « Modelizing Story » pp. 259-381. Undated; the PS gives
% « vers le six août ». The other pages of the batch (handwritten notes, listings on the
% backs) are left out, hence the gap in the page numbers. The leaf is scanned upside down.
%
% Private postal addresses are withheld (project rule): his own domicile after « Les
% Aumettes », and the street address of the second recipient. The first recipient's address
% is a university's and is kept.
\documentclass[11pt,a4paper]{article}
\input{../preamble/grothendieck}

\folder{134-8}
\batch{4}
\pages{61}{80}
\foldertitle{[Chapitre VII (pages 555 à 593) : Linearization of homotopy types] : tapuscrit annoté (22/10-12/11/[1983]), notes manuscrites (s.d.).}
\dating{[1983]}
\watermark{Édition de démonstration}

\begin{document}

\page{63}
\note{instructions dactylographiées, signées, pour la reproduction d'une partie de \emph{Pursuing Stacks} ; seule page du lot transcrite ici}

Recommandations pour la \struck{photocopie} \add{xérographie} de "Modelizing Story", pages 259 à
381 :

1) Attention \struck{aux} au centrage, pour ne pas couper une partie du texte

2) Faire attention aux pages 274 et 275,* qui par erreur ont été
frappées recto-verso (ne pas oublier donc la page 275 au verso !)
Faire une copie supplémentaire de la page 275.

3) \struck{\ill{}} \add{Faire la} copie en trois exemplaires, me renvoyer \struck{\ill{}} \add{l'original} à mon domicile
(Les Aumettes, \note{adresse privée omise}), et envoyer \struck{les} deux
\struck{\ill{}} \add{copies} aux adresses suivantes :

\begin{quote}
Ronnie Brown

School of Mathematics

University College of North Wales

Bangor, Gwynned LL 57 2 U W

Angleterre
\end{quote}

et

\begin{quote}
Zoghman Mebkhout

\note{adresse privée omise}

Paris
\end{quote}

Remettre la troisième copie à Monsieur Contou-Carrère.

\add{Avec mes remerciements}

\add{A Grothendieck}
\note{formule finale et signature autographes}

* Plus trois pages de table des matières.

PS Comme je suis absent de chez moi jusque vers le six août, il serait
préférable de ne renvoyer l'original que vers le quatre ou le cinq \uncertain{aoùt}\note{sic}.

\end{document}
